\chapter{Colposcopy}
\section*{What Is This Test?}
Colposcopy is a detailed examination of the cervix, vagina, and sometimes vulva using a magnifying instrument. Biopsy or endocervical sampling may be performed during the examination.
\section*{Alternative Names}
\begin{itemize}\item Colposcopic Examination\item Cervical Colposcopy\item Colposcopy with Biopsy\end{itemize}
\section*{Principle}
Acetic acid and sometimes iodine highlight abnormal epithelium. Magnified inspection guides biopsies of suspicious areas and helps determine whether cellular changes require observation, treatment, or further evaluation.
\section*{Why This Test Is Ordered}
It is commonly performed after an abnormal cervical screening result, persistent high-risk HPV, a suspicious cervical appearance, or selected unexplained bleeding.
\section*{Primary vs Secondary Test}
Colposcopy is a secondary diagnostic examination after cervical screening or clinical assessment. Histopathology from a directed biopsy provides the tissue diagnosis.
\section*{How This Test Is Performed}
With a speculum in place, the cervix is viewed through the colposcope, solutions are applied, and biopsies or an endocervical sample may be collected. The instrument does not enter the body.
\section*{What Preparation Is Needed}
Follow clinic instructions about vaginal intercourse, douching, and vaginal medicines. Tell the clinician about pregnancy possibility, anticoagulants, bleeding disorders, and allergies.
\section*{Contraindications and Precautions}
Cramping, spotting, and discharge are common after biopsy. Heavy bleeding, fever, severe pain, or foul discharge requires medical advice. Pregnancy changes which procedures are appropriate.
\section*{Understanding Results}
Pathology may show no dysplasia, low- or high-grade intraepithelial lesion, glandular abnormality, or invasive cancer. Management depends on histology, age, pregnancy status, and prior history.
\section*{Follow-up Tests}
Follow-up may include repeat HPV or cytology, excisional treatment such as LEEP, endocervical sampling, or gynecologic oncology referral.
\section*{References}
\begin{enumerate}\item MedlinePlus. Colposcopy.\item American Society for Colposcopy and Cervical Pathology. Management guidelines.\end{enumerate}
\clearpage\section*{Understanding Results and Follow-up}
The laboratory report should identify the specimen, method, units, reference interval or decision limit, and whether the result is positive, negative, abnormal, or inconclusive. Reference intervals vary with age, sex, pregnancy, specimen type, assay, and laboratory; a result outside a range is not automatically a diagnosis, and a result within a range does not always exclude disease. Discuss unexpected results with the ordering clinician, who can decide whether repeat or confirmatory testing, treatment, or referral is appropriate. Do not change medicines or delay urgent care based on an isolated result.

\section*{Regulatory Status and Authorization Date}This chapter describes a test category, not a specific commercial product. FDA, CDSCO, EU IVDR, and MHRA approval or registration dates therefore cannot be assigned without the assay manufacturer, product name, instrument, and intended use. For laboratory-developed tests, an individual agency approval date may not exist. See the Preface for the interpretation of regulatory status.
\clearpage
